\documentclass[11pt]{article}

\usepackage[T1]{fontenc}
\usepackage{lmodern}
\usepackage{microtype}
\usepackage{amsmath,amssymb,amsthm,mathtools}
\usepackage[a4paper,margin=25mm]{geometry}
\usepackage[dvipsnames]{xcolor}
\usepackage{fancyhdr}
\usepackage[hidelinks]{hyperref}

\definecolor{ink}{HTML}{1F2933}
\definecolor{accent}{HTML}{8C2F39}
\definecolor{statusbg}{HTML}{FFF0C9}
\definecolor{statusborder}{HTML}{8A5A00}

\hypersetup{
  pdftitle={PA00FG: The Gauss--Eisenstein Data Lemma},
  pdfsubject={A standard mathematical restatement of a native Peano Lab theorem},
  pdfauthor={Peano Lab arithmetic library},
  colorlinks=true,
  linkcolor=accent,
  urlcolor=accent
}

\pagestyle{fancy}
\fancyhf{}
\fancyhead[L]{\small\textcolor{accent}{PA00FG}}
\fancyhead[R]{\small Gauss--Eisenstein data lemma}
\fancyfoot[C]{\small\thepage}
\renewcommand{\headrulewidth}{0.4pt}

\setlength{\parindent}{0pt}
\setlength{\parskip}{0.2em}

\newtheorem{theorem}{Theorem}
\newcommand{\N}{\mathbb{N}}
\newcommand{\PrimePA}{\operatorname{Prime}_{\mathrm{PA}}}
\newcommand{\QRes}{\operatorname{QRes}}
\newcommand{\Even}{\operatorname{Even}}
\newcommand{\Odd}{\operatorname{Odd}}
\newcommand{\legendre}[2]{\left(\frac{#1}{#2}\right)}

\newcommand{\statusbox}[1]{%
  \begin{center}
    \fcolorbox{statusborder}{statusbg}{%
      \parbox{0.89\textwidth}{\small\textcolor{ink}{#1}}}
  \end{center}}

\begin{document}

\begin{center}
  {\color{accent}\Large\bfseries PA00FG}\par
  \vspace{0.35em}
  {\huge\bfseries The Gauss--Eisenstein Data Lemma}\par
  \vspace{0.55em}
  {\large A standard mathematical restatement of}\par
  \vspace{0.15em}
  {\large\ttfamily distinct\_odd\_primes\_gauss\_eisenstein\_data\_exists}
\end{center}

\statusbox{\textbf{Formal status.}
The native tactic body is kernel-checked relative to its declared dependencies,
but the theorem is currently a \texttt{candidate\_body\_checked} result rather
than a publicly admitted closed theorem.}

\section{Abbreviations for the expanded PA formulas}

All variables range over \(\N\).  The native theorem uses only \(0\),
successor, addition, multiplication, equality, first-order connectives, and
quantifiers.  The familiar notation below abbreviates formulas that are fully
expanded on the theorem page.

The PA primality predicate used in the statement is
\begin{equation}
  \PrimePA(r)
  \;:\Longleftrightarrow\;
  r\neq 1
  \;\land\;
  \forall a,b\in\N,
  \bigl(r=ab\Longrightarrow a=1\lor b=1\bigr).
  \label{eq:prime}
\end{equation}

Quadratic residuosity is encoded without subtraction:
\begin{equation}
  \QRes_m(a)
  \;:\Longleftrightarrow\;
  \exists x,u,v\in\N,
  \qquad x^2+mu=a+mv.
  \label{eq:qres}
\end{equation}
Thus \(\QRes_m(a)\) is the natural-number formulation of
\(x^2\equiv a\pmod m\).

Parity is expressed by
\begin{align}
  \Even(n)&\;:\Longleftrightarrow\;\exists t\in\N,\quad n=2t,
  \label{eq:even}\\
  \Odd(n)&\;:\Longleftrightarrow\;\exists t\in\N,\quad n=2t+1.
  \label{eq:odd}
\end{align}
Finally, balanced congruence modulo two means
\begin{equation}
  a\equiv b\pmod 2
  \;:\Longleftrightarrow\;
  \exists u,v\in\N,
  \qquad a+2u=b+2v.
  \label{eq:modtwo}
\end{equation}

\section{The clean mathematical statement}

\begin{theorem}[PA00FG]
Let \(p,q,h,k\in\N\).  Suppose
\begin{equation}
  p=2h+1,
  \qquad
  q=2k+1,
  \qquad
  \PrimePA(p),
  \qquad
  \PrimePA(q),
  \qquad
  p\neq q.
  \label{eq:hypotheses}
\end{equation}
Then there exist \(e,f,Q,U\in\N\) such that
\begin{align}
  \QRes_p(q)
    &\Longleftrightarrow \Even(e),
    \label{eq:first-residue}\\
  \neg\QRes_p(q)
    &\Longleftrightarrow \Odd(e),
    \label{eq:first-nonresidue}\\
  \QRes_q(p)
    &\Longleftrightarrow \Even(f),
    \label{eq:second-residue}\\
  \neg\QRes_q(p)
    &\Longleftrightarrow \Odd(f),
    \label{eq:second-nonresidue}\\
  e&\equiv Q\pmod 2,
    \label{eq:first-parity}\\
  f&\equiv U\pmod 2,
    \label{eq:second-parity}\\
  Q+U&=hk.
    \label{eq:eisenstein}
\end{align}
\end{theorem}

Because \(p=2h+1\) and \(q=2k+1\), the parameters are
\begin{equation}
  h=\frac{p-1}{2},
  \qquad
  k=\frac{q-1}{2}.
  \label{eq:halves}
\end{equation}
The equations in \eqref{eq:hypotheses} are used instead of subtraction or
division because neither operation is primitive in the native PA term
language.

\section{What the four witnesses represent}

The proof constructs \(e\) and \(f\) as the two Gauss-lemma sign counts.  In
ordinary notation, one may view them as
\begin{align}
  e
  &=\#\left\{
      1\leq i\leq h:
      \langle iq\rangle_p>\frac{p}{2}
    \right\},
  \label{eq:e-count}\\
  f
  &=\#\left\{
      1\leq j\leq k:
      \langle jp\rangle_q>\frac{q}{2}
    \right\},
  \label{eq:f-count}
\end{align}
where \(\langle n\rangle_m\) denotes the least positive residue of \(n\)
modulo \(m\).  Gauss's lemma explains
\eqref{eq:first-residue}--\eqref{eq:second-nonresidue}: quadratic residuosity
is exactly the evenness of the appropriate sign count.

The proof chooses \(Q\) and \(U\) as the Eisenstein quotient sums
\begin{align}
  Q&=\sum_{i=1}^{h}\left\lfloor\frac{iq}{p}\right\rfloor,
  \label{eq:q-sum}\\
  U&=\sum_{j=1}^{k}\left\lfloor\frac{jp}{q}\right\rfloor.
  \label{eq:u-sum}
\end{align}
The parity bridges \eqref{eq:first-parity} and
\eqref{eq:second-parity} identify the Gauss counts with these quotient sums
modulo two.  Eisenstein's lattice-point argument gives the exact identity
\begin{equation}
  \sum_{i=1}^{h}\left\lfloor\frac{iq}{p}\right\rfloor
  +
  \sum_{j=1}^{k}\left\lfloor\frac{jp}{q}\right\rfloor
  =hk.
  \label{eq:lattice}
\end{equation}

\paragraph{Formal nuance.}
The public surface of PA00FG exposes only the relations
\eqref{eq:first-residue}--\eqref{eq:eisenstein}.  The much larger internal
proof constructs the beta-coded finite sequences and sums that realize
\eqref{eq:e-count}--\eqref{eq:u-sum}, then hides those implementation
witnesses before returning \(e,f,Q,U\).

\section{Why this is the central quadratic-reciprocity package}

In conventional Legendre-symbol notation, the four classification formulas
say
\begin{equation}
  \legendre{q}{p}=(-1)^e,
  \qquad
  \legendre{p}{q}=(-1)^f.
  \label{eq:gauss-signs}
\end{equation}
The two parity bridges and the Eisenstein identity therefore give
\begin{align}
  \legendre{q}{p}\legendre{p}{q}
    &=(-1)^{e+f}
      \notag\\
    &=(-1)^{Q+U}
      \notag\\
    &=(-1)^{hk}
      \notag\\
    &=(-1)^{\frac{(p-1)(q-1)}{4}}.
  \label{eq:qr}
\end{align}
Equation \eqref{eq:qr} is the standard parity identity behind quadratic
reciprocity.  PA00FG is therefore not yet the final reciprocity theorem; it is
the common witness package from which the final same-status and
opposite-status cases are derived constructively.

\end{document}
